\documentclass[11pt]{article}
\usepackage[utf8]{inputenc}
\usepackage[T1]{fontenc}
\usepackage{geometry}
\usepackage{amsmath}
\usepackage{enumerate}
\geometry{a4paper, margin=2cm}

\begin{document}

\begin{center}
    {\Large \textbf{Gabarito -- PA e PG Matemática 1 Ano }} \\
    \vspace{0.3cm}
\end{center}

\begin{enumerate}

\item \textbf{Numa P.A., $a_1=3$ e $r=5$. Qual é o $6^\circ$ termo?}

\[
a_6 = a_1 + (6-1)r = 3 + 5\cdot5 = 28.
\]
\textbf{Resposta: D) 28}

\item \textbf{Termo geral da P.A. $(2,5,8,11,\ldots)$:}

\[
a_n = 2 + (n-1)\cdot3 = 3n - 1.
\]
\textbf{Resposta: } \(a_n=3n-1\)

\item \textbf{Soma dos 6 primeiros termos da P.A. $(3,8,13,\ldots)$:}

\[
a_1=3, \quad r=5,\quad a_6=3+5\cdot5=28,
\]
\[
S_6=\frac{6}{2}(a_1+a_6)=3(3+28)=93.
\]
\textbf{Resposta: } \(93\)

\item \textbf{Que número deve ser somado a cada termo de $(5,8,11)$ para obter $(9,12,15)$?}

\[
9 - 5 = 4
\]
\textbf{Resposta: C) 4}

\item \textbf{Qual a razão da P.A. $(-6,-3,0,3,6,\ldots)$?}

\[
r = -3 - (-6) = 3.
\]
\textbf{Resposta: C) 3}

\item \textbf{O $5^\circ$ termo de uma P.G. com $a_1=2$ e $q=3$:}

\[
a_5 = 2 \cdot 3^{4} = 2 \cdot 81 = 162.
\]
\textbf{Resposta: D) 162}

\item \textbf{Qual é a razão da P.G. $(2,6,18,54,\ldots)$?}

\[
q=\frac{6}{2}=3.
\]
\textbf{Resposta: B) 3}

\item \textbf{Em uma P.G. de razão $\frac12$ e $a_1=64$, o $4^\circ$ termo é:}

\[
a_4 = 64\left(\tfrac12\right)^3 = 64 \cdot \tfrac18 = 8.
\]
\textbf{Resposta: C) 8}

\item \textbf{A P.G. $(3,6,12,\ldots)$ é crescente, decrescente ou constante? Justifique.}

\[
q = \frac{6}{3} = 2 > 1.
\]
Logo, cada termo é maior que o anterior.  
\textbf{Resposta: Crescente.}

\item \textbf{O produto dos três primeiros termos de uma P.G. é 8, e a razão é 2. O primeiro termo vale:}

Se \(a_1 = x\), então:
\[
x \cdot 2x \cdot 4x = 8x^3 = 8
\]
\[
x^3 = 1 \Rightarrow x = 1.
\]
\textbf{Resposta: B) 1}

\end{enumerate}

\end{document}

